\documentclass[a4paper, serif, 10pt]{moderncv}

%% moderncv
% style and color
\moderncvstyle{banking}
\moderncvcolor{black}

%% page layout/margins
\usepackage[top=2.5cm, bottom=2.5cm, left=1.5cm, right=1.5cm]{geometry}

\nopagenumbers{}

%% Babel config for Indonesian (Bahasa Indonesia) language
% ref:
% - https://latex3.github.io/babel/guides/locale-indonesian.html
\usepackage[indonesian]{babel}

%% fontspec
\usepackage{fontspec}

\IfFontExistsTF{Linux Libertine}{
  \setmainfont[Ligatures={TeX, Common}, Numbers=OldStyle]{Linux Libertine}
}{}
\IfFontExistsTF{Linux Libertine O}{
  \setmainfont[Ligatures={TeX, Common}, Numbers=OldStyle]{Linux Libertine O}
}{}

%% CTeX
% CJK support, used to show CJK characters in the resume
%
% - fontset=none: disable builtin fontset but instead set the CJK font manually
% - heading=false: disable ctex heading
% - punct=kaiming: use kaiming punctuations styles for CJK
% - scheme=plain: use plain scheme, do not override \normalsize font size
% - space=auto: space settings for CJK characters
%
% ref:
% - http://ctan.mirrorcatalogs.com/language/chinese/ctex/ctex.pdf
\usepackage[UTF8, heading=false, punct=kaiming, scheme=plain, space=auto]{ctex}

\IfFontExistsTF{Noto Serif CJK SC}{
  \setCJKmainfont{Noto Serif CJK SC}
}{}
\IfFontExistsTF{Noto Sans CJK SC}{
  \setCJKsansfont{Noto Sans CJK SC}
}{}

%% line spacing
\usepackage{setspace}
\setstretch{1.125}

%% URL styling - use normal text instead of monospace
\urlstyle{same}

% Auto-underline all links
% Defer until \begin{document} so hyperref is already loaded
\AtBeginDocument{%
  \let\oldhref\href
  \renewcommand{\href}[2]{\oldhref{#1}{\underline{#2}}}%
}

%% Patch \httplink and \httpslink to handle full URLs
% The original moderncv macros always prepend http:// or https:// to URLs.
% This causes issues when the URL already has a protocol (e.g., https://example.com)
% resulting in malformed URLs like https://https://example.com.
% This patch checks if the URL already contains "://" and uses it directly if so.
% Note: We use \str_set:Nx (x-expansion) to expand macros like \@homepage first.
\ExplSyntaxOn
\str_new:N \g_yamlresume_url_str

\RenewDocumentCommand{\httpslink}{O{}m}{
  \str_set:Nx \g_yamlresume_url_str {#2}
  \str_if_in:NnTF \g_yamlresume_url_str {://}
    { \url{#2} }
    { \url{https://#2} }
}

\RenewDocumentCommand{\httplink}{O{}m}{
  \str_set:Nx \g_yamlresume_url_str {#2}
  \str_if_in:NnTF \g_yamlresume_url_str {://}
    { \url{#2} }
    { \url{http://#2} }
}
\ExplSyntaxOff

\name{Budi Santoso}{}
\title{Insinyur DevOps | Otomasi Infrastruktur \& Cloud}
\phone[mobile]{+62 812-3456-7890}
\email{budi.santoso.devops@example.com}
\homepage{https://budisantoso.dev}

\address{Jl. Jenderal Sudirman No. 52, Jakarta Pusat, Jakarta, DKI Jakarta, Indonesia, 10220}{}{}

\extrainfo{{\small \faGithub}\ \href{https://github.com/budisantoso-devops}{@budisantoso-devops} {} {} {} • {} {} {} 
{\small \faLinkedin}\ \href{https://www.linkedin.com/in/budisantoso-devops}{@budisantoso-devops}}

\begin{document}

\maketitle

\section{Informasi Dasar}

\cvline{}{\begin{itemize}
\item Insinyur DevOps dengan lebih dari 6 tahun pengalaman dalam merancang, mengotomasi, dan mengoperasikan infrastruktur cloud yang andal.
\item Terampil dalam Kubernetes, Terraform, CI/CD, dan praktik Site Reliability Engineering untuk meningkatkan kecepatan rilis serta ketersediaan layanan.
\item Berpengalaman membangun platform observabilitas, mengelola multi-cloud, dan menerapkan keamanan pada pipeline DevOps.
\item Kolaboratif, berorientasi pada otomasi, dan terbiasa memimpin transformasi budaya DevOps di tim lintas fungsi.
\end{itemize}}
\section{Pendidikan}

\cventry{Agu 2015–Jul 2019}
        {Sarjana, Teknik Informatika, Nilai: 3.72}
        {Universitas Indonesia}
        {\url{https://www.ui.ac.id}}
        {}
        {\begin{itemize}
\item Lulus dengan predikat cum laude dan fokus pada rekayasa perangkat lunak serta infrastruktur.
\item Mengembangkan proyek akhir tentang orkestrasi kontainer untuk aplikasi berbasis mikroservis.
\item Aktif dalam organisasi kemahasiswaan yang bergerak di bidang teknologi dan open source.
\end{itemize}
\textbf{Mata Kuliah}: Algoritma dan Struktur Data, Sistem Operasi, Jaringan Komputer, Basis Data, Arsitektur Perangkat Lunak, Keamanan Informasi, Komputasi Awan, Otomasi Infrastruktur}
\section{Pengalaman Kerja}

\cventry{Jan 2022–Sekarang}
        {Insinyur DevOps Senior}
        {PT Solusi Digital Nusantara}
        {\url{https://www.solusidigitalnusantara.co.id}}
        {}
        {\begin{itemize}
\item Memimpin adopsi GitOps dan pipeline CI/CD untuk lebih dari 40 layanan mikro, menurunkan waktu deployment dari 2 jam menjadi 15 menit.
\item Mengelola cluster Kubernetes multi-tenant di AWS dengan Terraform, Helm, dan Argo CD untuk menjaga konsistensi lingkungan.
\item Membangun platform observabilitas terpusat menggunakan Prometheus, Grafana, Loki, dan OpenTelemetry.
\item Menerapkan kebijakan keamanan supply chain, pemindaian kerentanan, dan manajemen rahasia dengan Vault.
\item Menjadi mentor bagi 4 insinyur junior dan memimpin forum praktik terbaik DevOps internal.
\end{itemize}
\textbf{Kata Kunci}: Kubernetes, Terraform, AWS, GitLab CI, Observabilitas}

\cventry{Mar 2020–Des 2021}
        {Insinyur DevOps}
        {PT Teknologi Cakra}
        {\url{https://www.teknologicakra.co.id}}
        {}
        {\begin{itemize}
\item Mengotomasi provisioning infrastruktur dengan Ansible dan Terraform, mengurangi kesalahan konfigurasi manual hingga 70\%.
\item Membangun pipeline Jenkins untuk build, test, dan deployment aplikasi Java dan Node.js.
\item Mengimplementasikan monitoring dan alerting dengan Prometheus, Alertmanager, dan Grafana.
\item Melakukan tuning Linux, Docker, dan Nginx untuk meningkatkan performa serta stabilitas layanan.
\end{itemize}
\textbf{Kata Kunci}: Docker, Jenkins, Ansible, Prometheus, Linux}

\cventry{Jun 2019–Feb 2020}
        {Site Reliability Engineer}
        {PT Data Prima}
        {\url{https://www.dataprima.co.id}}
        {}
        {\begin{itemize}
\item Menjaga ketersediaan layanan data dengan pemantauan proaktif dan respons insiden.
\item Menulis skrip otomasi Python dan Bash untuk tugas operasional berulang.
\item Mendokumentasikan runbook, prosedur pemulihan bencana, dan praktik on-call.
\end{itemize}
\textbf{Kata Kunci}: Python, Bash, Nagios, MySQL, Jaringan}
\section{Bahasa}

\cvline{Melayu}{Bahasa Ibu / Bilingual \hfill \textbf{Kata Kunci}: Bahasa Indonesia}
\cvline{Inggris}{Tingkat Profesional Penuh \hfill \textbf{Kata Kunci}: IELTS 7.0, TOEFL 600}
\section{Keahlian}

\cvline{Orkestrasi Kontainer}{Ahli \hfill \textbf{Kata Kunci}: Docker, Kubernetes, Helm, OpenShift}
\cvline{Infrastruktur sebagai Kode}{Mahir \hfill \textbf{Kata Kunci}: Terraform, Ansible, Pulumi, CloudFormation}
\cvline{CI/CD}{Ahli \hfill \textbf{Kata Kunci}: GitLab CI, Jenkins, GitHub Actions, Argo CD}
\cvline{Cloud}{Mahir \hfill \textbf{Kata Kunci}: AWS, Google Cloud, Azure, Alibaba Cloud}
\cvline{Observabilitas}{Mahir \hfill \textbf{Kata Kunci}: Prometheus, Grafana, ELK Stack, Datadog}
\cvline{Bahasa Pemrograman}{Mahir \hfill \textbf{Kata Kunci}: Python, Go, Bash, SQL}
\cvline{Keamanan DevOps}{Menengah \hfill \textbf{Kata Kunci}: Vault, Trivy, SonarQube, IAM}
\section{Penghargaan}

\cventry{Des 2023}
        {PT Solusi Digital Nusantara}
        {Penghargaan Inovasi Otomasi Infrastruktur}
        {}
        {}
        {Diberikan atas kontribusi dalam membangun pipeline CI/CD terpusat yang meningkatkan kecepatan rilis dan keandalan layanan.}

\cventry{Agu 2021}
        {Cloud Native Indonesia}
        {Juara 2 Kompetisi Cloud Native}
        {}
        {}
        {Meraih juara kedua dalam kompetisi desain arsitektur cloud native untuk studi kasus layanan logistik.}
\section{Sertifikat}

\cventry{Mar 2022}
        {Amazon Web Services}
        {AWS Certified DevOps Engineer - Professional}
        {\url{https://aws.amazon.com/certification/certified-devops-engineer-professional/}}
        {}
        {}

\cventry{Nov 2021}
        {Cloud Native Computing Foundation}
        {Certified Kubernetes Administrator}
        {\url{https://training.linuxfoundation.org/certification/certified-kubernetes-administrator-cka/}}
        {}
        {}

\cventry{Jun 2023}
        {HashiCorp}
        {HashiCorp Certified: Terraform Associate}
        {\url{https://www.hashicorp.com/certification/terraform-associate}}
        {}
        {}
\section{Publikasi}

\cventry{Agu 2023}
        {Otomasi Infrastruktur Multi-Cloud dengan Terraform dan GitLab CI}
        {Jurnal Teknologi Cloud Indonesia}
        {\url{https://jurnalcloud.id/artikel/otomasi-infrastruktur-multi-cloud}}
        {}
        {\begin{itemize}
\item Membahas pola modular Terraform untuk mengelola sumber daya di AWS, GCP, dan Azure.
\item Menjelaskan integrasi pipeline GitLab CI untuk deployment yang aman dan dapat diaudit.
\end{itemize}}

\cventry{Okt 2022}
        {Observabilitas Modern untuk Arsitektur Microservices}
        {Konferensi DevOps Indonesia}
        {\url{https://devops.id/prosiding/observabilitas-microservices}}
        {}
        {\begin{itemize}
\item Mengulas penerapan metrik, log, dan trace terdistribusi pada layanan mikro.
\item Memberikan rekomendasi praktis untuk alerting dan penanganan insiden.
\end{itemize}}
\section{Proyek}

\cventry{Jan 2023–Jun 2023}
        {Implementasi observabilitas terpusat untuk 50+ layanan mikro.}
        {Platform Observabilitas Terpusat}
        {\url{https://github.com/budisantoso-devops/observability-platform}}
        {}
        {\begin{itemize}
\item Menggabungkan metrik, log, dan trace menggunakan Prometheus, Loki, dan OpenTelemetry.
\item Membuat dasbor Grafana dan alerting yang mengurangi mean time to recovery sebesar 40\%.
\end{itemize}
\textbf{Kata Kunci}: Prometheus, Grafana, Loki, OpenTelemetry}

\cventry{Mar 2022–Nov 2022}
        {Pipeline otomatis untuk deployment ke AWS, GCP, dan Azure.}
        {Pipeline CI/CD Multi-Cloud}
        {\url{https://github.com/budisantoso-devops/multi-cloud-cicd}}
        {}
        {\begin{itemize}
\item Membangun pipeline reusable dengan GitLab CI, Terraform, dan Argo CD.
\item Menerapkan quality gate, pemindaian keamanan, dan strategi canary release.
\end{itemize}
\textbf{Kata Kunci}: GitLab CI, Terraform, Argo CD, Kubernetes}
\section{Minat}

\cvline{Olahraga}{Lari, Bersepeda, Renang}
\cvline{Musik}{Gitar, Piano}
\cvline{Open Source}{Kubernetes, Terraform, Prometheus}
\section{Kegiatan Sukarela}

\cventry{Jan 2022–Sekarang}
        {Mentor}
        {Komunitas DevOps Indonesia}
        {\url{https://devops.id}}
        {}
        {\begin{itemize}
\item Menjadi mentor bagi anggota komunitas dalam mempelajari Kubernetes, Terraform, dan praktik CI/CD.
\item Membantu menyelenggarakan workshop dan sesi diskusi teknis rutin.
\end{itemize}}
    

\cventry{Sep 2017–Jul 2019}
        {Pengajar}
        {Universitas Indonesia Programming Club}
        {\url{https://ui.ac.id/programming-club}}
        {}
        {\begin{itemize}
\item Mengajarkan dasar pemrograman dan penggunaan Git kepada mahasiswa baru.
\item Mendampingi peserta dalam kompetisi pemrograman tingkat universitas.
\end{itemize}}
    
\end{document}